Cieľom tejto bakalárskej práce bolo navrhnúť a implementovať malé vozidlo schopné autonómneho riadenia a detekcie prekážok s využitím metód počítačového videnia. Vozidlo bolo vybavené kamerovým modulom ESP32-CAM a ultrazvukovým senzorom, ktoré nepretržite monitorovali priestor pred sebou a poskytovali dáta pre riadiaci server pracujúci v programovacom jazyku Python. Dôležitým prvkom riešenia bolo spracovanie obrazu v knižnici OpenCV, kde sa prostredníctvom rozlíšenia farieb a hľadania kontúr v definovanej oblasti záujmu zabezpečovalo udržiavanie vozidla v stredovej línii vozovky. Na detekciu prekážok bola do systému integrovaná konvolučná neurónová sieť YOLOv8, ktorá spolu s meraním vzdialenosti umožnila vozidlu plynulo reagovať na objekty. Celý hardvérový systém bol riadený platformou Arduino, pričom spracovanie obrazu prebiehalo na serveri, ktorý bezdrôtovo pomocou Wi-Fi komunikoval s vozidlom.
